% resumo em inglês
\setlength{\absparsep}{18pt} % ajusta o espaçamento dos parágrafos do resumo
\begin{resumo}[Abstract]
 \begin{otherlanguage*}{english}
<<<<<<< HEAD
    The adoption of Building Information Modeling (BIM) has been digitally transforming the construction industry by centralizing information throughout the entire project lifecycle, increasing accuracy, reducing errors, and optimizing resources. This growth, however, brings challenges related to the management of large volumes of data and, above all, to the integration of complex regulatory standards. In the Architecture, Engineering, and Construction (AEC) sector, automatic compliance checking is still hindered by technical language and the lack of standardization of normative documents. A promising alternative is the RASE methodology (\textit{Requirement, Applicability, Selection, Exception}), which converts standards into computable data. In parallel, Large Language Models (LLMs) have been investigated to interpret regulatory texts, but systematic evaluations of open-source LLMs applied to Brazilian standards are still scarce. This research comparatively evaluates six open-source LLMs --- Llama, Dolphin, Gemma, Mistral, Alpaca, and Qwen --- for the automatic conversion of AEC technical standards into the RASE format, using \textbf{exclusively Prompt Engineering} (without Fine-Tuning and without Retrieval-Augmented Generation). The task is decomposed into three successive normalization levels: N1 --- atomic segmentation of normative sentences; N2 --- identification of the RASE operators (Requirement, Applicability, Selection, Exception); and N3 --- semantic structuring into JSON. Three experiments were conducted: EN1 (N1 in isolation), EN2 (N2 in isolation, from the reference N1), and EN1N2 (chained pipeline), allowing the analysis of error propagation between stages. The models were served locally via Ollama and applied to a dataset of 79 annotated Brazilian standards, based on NBR 9050. Output quality was assessed by thirteen metrics, organized in five families: lexical similarity (FuzzyWuzzy, TF-IDF), semantic similarity (SBERT, BERTimbau, Multilingual), distance metrics (Word Mover's Distance with FastText and NILC embeddings), generation-oriented metrics (BERTScore and ROUGE-L), and classification metrics (Accuracy, Precision, Recall, and F1-score). Considering the global average across the nine similarity metrics over the three experiments, Dolphin leads (0.676), followed by Llama (0.658) and Mistral (0.654); Dolphin also offers the best cost-benefit, with execution time about ninety times lower than that of Llama in the chained pipeline. Llama leads the EN2 experiment and the isolated N3 level in semantic metrics. Semantic metrics and BERTScore proved more suitable than lexical metrics for evaluating normative transformations, with BERTimbau showing a consistent gain over SBERT in Brazilian Portuguese. The classification metrics revealed that the Selection and Exception operators are the hardest classes in N2, with F1 close to zero for all models. The research contributes with the operational definition of the three normalization levels over the RASE methodology, with a reproducible set of prompts specialized by level and operator, and with the first systematic comparison of six open-source LLMs for the conversion of Brazilian AEC standards.
=======
    The adoption of Building Information Modeling (BIM) has been digitally transforming the construction industry by centralizing information throughout the entire project lifecycle, increasing accuracy, reducing errors, and optimizing resources. This growth, however, brings challenges related to the management of large volumes of data and, above all, to the integration of complex regulatory standards. In the Architecture, Engineering, and Construction (AEC) sector, automatic compliance checking is still hindered by technical language and the lack of standardization of normative documents. A promising alternative is the RASE methodology (\textit{Requirement, Applicability, Selection, Exception}), which converts standards into computable data. In parallel, Large Language Models (LLMs) have been investigated to interpret regulatory texts, but systematic evaluations of open-source LLMs applied to Brazilian standards are still scarce. This research comparatively evaluates six open-source LLMs --- Llama, Dolphin, Gemma, Mistral, Alpaca, and Qwen --- for the automatic conversion of AEC technical standards into the RASE format, using \textbf{Prompt Engineering}. The task is decomposed into three successive normalization levels: N1 --- atomic segmentation of normative sentences; N2 --- identification of the RASE operators (Requirement, Applicability, Selection, Exception); and N3 --- semantic structuring into JSON. Three experiments were conducted: EN1 (N1 in isolation), EN2 (N2 in isolation, from the reference N1), and EN1N2 (chained pipeline), allowing the analysis of error propagation between stages. The models were served locally via Ollama and applied to a dataset of 79 annotated Brazilian standards, based on NBR 9050. Output quality was assessed by lexical similarity metrics (FuzzyWuzzy, TF-IDF), semantic similarity metrics (SBERT, BERTimbau, Multilingual), and distance metrics (Word Mover's Distance with FastText and NILC embeddings), complemented by classification metrics (Accuracy, Precision, Recall, and F1-score). Results indicate that Llama leads in overall quality, achieving semantic similarity above 0.80 in the EN2 experiment, while Dolphin shows the best cost-benefit by combining competitive quality with execution time up to one hundred times lower. Semantic metrics proved to be more suitable than lexical ones for evaluating normative transformations, and the chained pipeline highlighted the expected propagation of errors between levels. The research contributes with the operational definition of the three normalization levels over the RASE methodology, with a reproducible set of prompts specialized by level and operator, and with the first systematic comparison of six open-source LLMs for the conversion of Brazilian AEC standards.
>>>>>>> b1bd44c (citation)
   \vspace{\onelineskip}

   \noindent
   \textbf{Keywords}: Machine Learning, Natural Language Processing, LLM, Prompt Engineering, BIM, RASE.
 \end{otherlanguage*}
\end{resumo}
